% ===== §8 =====
\section{Entry into the Symbolic and the Pre-Symbolic Remainder}
\label{p26:sec:entry}

\noindent
This section opens Part Two by tracing the shared form of the six limits to its source. The aim is to establish that the remainder language could not reach, in each of the six dimensions, has a single origin in the constitution of the speaking subject, and that psychoanalysis names this origin: the subject is constituted in the signifier by an operation that leaves a part outside itself, in the register of the Real. The section sets out the psychoanalytic account of entry into language, follows the two moments of that entry by which the remainder is produced, distinguishes the sense of the Real it employs, and shows that the six cores of Part One are expressions of this one remainder.

\subsection*{From six limits to one source}

Part One left a result and a question. The result was that in each dimension of relational life, language reaches toward a core it does not reach, and that in each the not-reaching is generative. The question is why the same form should recur across dimensions as different as the existence of a bond, the measure of an ethical response, and the direction of growth. Six limits with one form ask for one source, and the source is not to be found by comparing the dimensions with one another but by descending beneath them, to the constitution of the subject who speaks in all of them.

\subsection*{The subject constituted in the signifier}

To become a subject who speaks is to be taken up into the symbolic order, the field of signifiers that precedes the subject and structures it. This is the pivot of Lacan's rereading of Freud, and its consequence is precise. The subject that comes to be in the signifier is not a prior fullness that language then expresses. It is constituted in and by the signifier, and there is no subject of speech behind or before this constitution to which one might appeal \parencite{lacan2006ecrits}. Fink states the mechanism carefully: the subject is what is represented by one signifier for another, a subject that appears only in the movement of the chain and has no substance apart from it \parencite{fink1995subject}. The subject is an effect of the signifier, not its origin.

Before the signifier there is a phase in which the human infant finds an anticipated unity in an external image, a coherence borrowed from a form outside itself, and this borrowed coherence already installs a distance between the subject and the ground of its own being \parencite{lacan2006ecrits}. The image gives a unity the body does not yet have, and the self that forms around it is from the first an identification with something other, a self founded on an outside. When the signifier arrives it does not heal this distance but redoubles it, for now the subject must find itself in a field of signs that likewise precede and exceed it. The entry into language is the second and decisive founding of the subject upon what is not itself.

\subsection*{The two moments of alienation and separation}

The entry into the signifier can be followed in two moments, and the second is where the remainder is produced. In the first moment, which psychoanalysis calls alienation, the subject takes on the signifiers of the Other as the only means of appearing at all, and pays for its appearance by a loss: to be represented by a signifier is to be reduced to what the signifier can carry, and the living being that entered is not wholly carried \parencite{lacan1977xi}. In the second moment, separation, the subject confronts a lack in the Other itself, the fact that the Other's chain of signifiers does not close, does not contain the answer to what the subject is for it, and in this confrontation the object-cause of desire is produced, at the point where the two lacks meet \parencite{fink1995subject}. What matters for the present argument is that this double operation does not gather the whole of the subject into the signifier. It constitutes the subject by a cut, and the cut leaves a portion of the subject's being outside the signifier, uncaptured by the structure that made the subject a subject.

Freud had already found, beneath the subject of conscious speech, a domain governed otherwise, and had named the pressure that rises from it the drive, \emph{Trieb}, a force on the border between the somatic and the psychical that presses toward satisfaction and is never simply stilled \parencite{freud1915instincts}. Lacan's contribution was to show that the induction into the signifier is what produces this border, that the very operation constituting the subject as a speaker leaves the drive's portion outside the signifier. This uncaptured portion is the remainder, and it is not a residue that a fuller symbolization might later gather in. It is what the operation of symbolization, by its nature, leaves out.

\begin{casebox}{the word that never fit}
A person carries, from early life, a way of being wanted that no word they were given ever fit, a pressure with no name in the language handed to them. In adult intimacy it presses still, shaping what they seek and cannot say they seek. It is not that the right word was withheld and might yet be supplied. There was never a word of the right kind, because what presses is exactly the portion the entry into words left out. The nameless pressure is not a failure of their vocabulary. It is the remainder, felt.
\end{casebox}

\subsection*{Which Real}

The word Real must be used with care, since Lacan gave it more than one sense across his teaching. In the early work the Real is close to the brute given, what is always in its place, prior to the symbolic articulation that will carve it \parencite{lacan1993iii}. In the mature and late work it becomes the impossible, the internal limit of symbolization, what does not cease not to be written \parencite{lacan1998xx}. The present paper takes the Real primarily in the second, mature sense, as the internal limit rather than an external given, and it takes up the pre-symbolic remainder as the genetic aspect of that same limit, not as a separate brute domain lying behind the sign. The next section shows that these are two faces of one gap. Here the point is only to fix the usage: the Real that concerns the paper is the remainder produced by the operation of the signifier, the internal offcut of subjectivation, and not a fullness lying behind language that language fails to reach.

\subsection*{The remainder as the source of the six cores}

The pre-symbolic remainder, understood this way, is the source the six limits pointed to. What language could not reach, in existence, in maintenance, in value, in ethics, in justice, in evolution, was in each case a form of this same remainder, the part that the operation of the signifier leaves outside itself. The existence of the bond, the living present, the relational generation of value, the aesthetic measure, the generativity justice guards, the direction of growth: each is a face of the one thing that the constitution of the speaking subject cannot take up into the signifier. The six dimensions were six windows onto a single gap. The next section gives the gap its two faces and shows that they are one, and the section after gives the gap its object.
